% =============================================================
% Introduction
% Chapter order is set in ../main.tex; the rendered chapter number
% comes from \chapter{} below.
% =============================================================

\chapter{Introduction}
\label{ch:intro}

For most of computing's history, ``artificial intelligence'' was a
stage prop rather than a tool. Generations of science fiction
trained us to expect AI to arrive as a singular event --- HAL
waking up, Skynet going live, Data being switched on, the doors of
the Enterprise's main computer parting for a voice that knew
everything. What actually happened, beginning in late~2022 and
accelerating without pause through~2026, was both less dramatic and
more disruptive than the fiction predicted. Language models stopped
being curiosities and started being useful. The change crept into
editors, into terminals, into spreadsheets, into the way people
read papers and wrote emails, before anyone had agreed on a
vocabulary for what was happening.

\noindent
This book is for the engineer asked to build something on top of
that change.

\section{Where We Are}
\label{sec:intro-now}

The most reliable way to take stock of what large language models
can do today is to look at what has already been done with them.
Three examples, each in a different domain, are enough to set the
scale.

In February~2026, Anthropic published the result of an internal
experiment in which sixteen instances of Claude Opus~4.6, working
in parallel and coordinating through a shared workspace, wrote a
hundred-thousand-line C compiler in Rust, with backends for x86,
ARM, and RISC-V, capable of building the Linux~6.9
kernel~\citep{anthropic2026compiler}. The agents ran for roughly
two thousand sessions and consumed around twenty thousand dollars
of API credit. There was no central architect; there was no human
reviewing each commit. The compiler that came out is not a drop-in
replacement for GCC, but it boots an operating system.

In December~2025, the Erd\H{o}s problem~\#728 --- a number-theoretic
conjecture that had been open for decades --- was settled by a pair
of language-model systems working in tandem: GPT-5.2~Pro for
informal reasoning and the Aristotle theorem prover for
\emph{Lean}-verified formal proof~\citep{tao2025erdos}. The
resolution is the first Erd\H{o}s problem widely regarded as having
been solved autonomously by an AI system. The verifier rules out
hallucinated proofs by construction, which is what makes the result
trustworthy in a way that an unaided model's output is not. By
May~2026 the pattern had been industrialized: a Google DeepMind
team built \emph{AlphaProof Nexus}, a framework in which Gemini~3.1
Pro subagents propose Lean proof sketches and the Lean compiler
verifies them, optionally coordinated by an evolutionary search;
their full-featured agent autonomously resolved 9 of 353 open
Erd\H{o}s problems and proved 44 of 492 open conjectures from the
Online Encyclopedia of Integer Sequences, at a few hundred dollars
per problem~\citep{tsoukalas2026alphaproofnexus}. The same paper
reports that a stripped-down agent --- just LLM-generation
alternating with Lean verification, no specialized prover ---
solved all 9 Erd\H{o}s problems too, though more expensively on
the hardest. As LLMs get stronger, the specialized scaffolding
shrinks; the loop does not.

In January~2026, a research group at KAIST reported a workflow in
which a language model proposes structural redesigns of materials
that are difficult to synthesize, with the proposals filtered by
chemistry-informed rules. Of the top hundred candidates, thirty-four
matched experimentally synthesized structures that were not in the
training data~\citep{kaist2026materials}. The model did not discover
new chemistry; it explored a vast combinatorial space too quickly
for a human to match, with the chemistry kept honest by tools the
model could call.

In May~2026, a separate result from OpenAI shifted the upper
bound of what the math community considers achievable: a
general-purpose reasoning model --- not a system trained for
mathematics, not scaffolded toward proof search, not targeted at
this problem --- \emph{disproved} Erd\H{o}s's planar
unit-distance conjecture~\citep{openai2026discrete}. The proof
imports infinite class field towers and Golod--Shafarevich theory
into a purely combinatorial-geometry question, a cross-domain
bridge nobody had expected. External mathematicians, led by Will
Sawin at Princeton, verified the construction and pinned down a
fixed $\delta = 0.014$ for the new $n^{1+\delta}$ lower bound;
Tim Gowers called the result ``a milestone in AI mathematics.''
Unlike the formal-proof results above, the OpenAI disproof did
not lean on a Lean verifier or any specialized harness; the
system around the model was a batch evaluation that surfaced the
output and a community of mathematicians who checked it after the
fact.

The common thread runs through every example, though to different
degrees. None of these is a naked language-model achievement;
each is a \emph{system} --- a model in a loop with some
combination of executors, verifiers, search procedures, and human
or machine oversight. The mix varies. The C compiler is mostly
model-plus-orchestration. The Erd\H{o}s and OEIS results are
mostly model-plus-formal-verifier. The KAIST materials work is
mostly model-plus-domain-rules. The discrete-geometry disproof is
mostly model-plus-external-review. But the pattern is consistent:
the model is the engine; the system is the vehicle.

\section{A Moving Target}
\label{sec:intro-target}

The texture of the field is one of constant motion. Frontier model
families are revised every few months. Per-token prices have fallen
by more than an order of magnitude since~2023. The lineup of major
providers has reshuffled twice in the time it has taken to write
this book. Open-weights models, scarcely usable for serious work in
2023, now run usefully on a developer's laptop, and several of them
sit close to the closed frontier. Specialized tooling --- agent
harnesses, vector stores, eval frameworks, integration protocols ---
churns at a rate that would embarrass any other corner of software
engineering.

Two consequences for the reader. First, anything written about
specific models, prices, or rankings is out of date by the time it
is printed. We will name names where it helps the reader place
themselves --- Anthropic's Claude, OpenAI's GPT, Google's Gemini,
Meta's Llama, Alibaba's Qwen, DeepSeek, Mistral, Google's Gemma ---
but we will not stake the argument on which model is best at what.
Second, the architectural patterns underneath the churn have proved
more stable than one might fear. The agent loop, the tool-call
abstraction, the chat-message structure, the trade-offs between
context, cost, and capability --- these were already legible
in~2024 and have not changed shape since. \emph{Those} are what we
write about. Where the book makes a recommendation, it is about a
pattern, not a product.

A note on tone. The hype around language models is real, loud, and
will not subside on its own. The dismissiveness around them, in
some quarters, is also real. Both attitudes lead the engineer
astray. The most useful disposition for the reader of this book is
the one a working programmer brings to any other powerful tool: it
does some things well, others poorly; it has failure modes that
matter; it rewards careful design and punishes lazy assumptions; it
is now part of the toolchain whether one likes it or not.

\section{How We Got Here}
\label{sec:intro-history}

A history of how language models became what they are could fill a
book of its own. Three milestones are enough to orient the reader.

The first is the architectural shift. In~2017,
\citet{vaswani2017attention} introduced the transformer: a neural
network architecture that processed sequences using attention
rather than recurrence, scaled almost embarrassingly well with
parameter count and data, and proved durable enough that almost
every language model of consequence since has been a variant of it.
The decade between the transformer paper and now is a story of
training larger and larger models on more and more data, and of
discovering, repeatedly, that capability emerges from scale in ways
the architects had not predicted.

The second is the public moment. In November~2022, OpenAI released
ChatGPT --- a chat interface to a fine-tuned variant of GPT-3.5 ---
and the world abruptly noticed. The model was not, by the technical
standards of the field, a leap; the interface was. For the first
time, a non-specialist could ask the system to do useful things in
plain English and get useful things back. The downstream effects
were both political (every large company suddenly had an AI
strategy) and technical (every other lab raced to ship competitors;
open-weights releases proliferated; price-per-token began its
collapse).

The third milestone is the one this book starts from. In~2023, the
agent turn began: \citet{yao2023react} formalized the
reasoning-and-action loop; OpenAI shipped function-calling APIs that
let models emit structured tool requests; \citet{wang2024codeact}
argued that the right action shape was executable code; Anthropic
introduced the Model Context Protocol in late~2024, standardizing
the way agents discover tools. Each of these moves added one piece
of the substrate on which today's agents are built. The C compiler,
the Erd\H{o}s proof, and the materials redesign of
\S\ref{sec:intro-now} are among the first systems to use the whole
substrate at once.

\section{From Models to Systems}
\label{sec:intro-systems}

If the reader takes one thing from this introduction, let it be
this: a contemporary language model is not, by itself, the artifact
worth designing around. The artifact is the \emph{system} that
wraps the model. Recognizing this changes the questions one asks at
the architecture stage, the metrics one tracks at the operations
stage, and the failures one prepares for at the production stage.

We will use the word \emph{agent} throughout this book in the
following narrow sense.

\begin{definitionbox}
An \textbf{agent} is a system in which a language model
repeatedly takes actions, observes the results of those actions,
and decides what to do next, until it judges the assigned task
complete. The actions are taken via \emph{tools}: external pieces
of code that the model can invoke and whose return values are fed
back into its context. The system is the assembly of model, tool
inventory, control loop, memory, and policy that turns a single
turn of generation into work in the world.
\end{definitionbox}

\noindent
Returning to the achievements of \S\ref{sec:intro-now} with this
definition in hand: the C compiler is an agent system in which
each Claude instance had a shell tool, a workspace, a build runner,
and a coordination protocol with its peers; the Erd\H{o}s result
is an agent system in which the language model proposed proof
sketches that a Lean verifier accepted or rejected; the materials
redesign is an agent system in which the model generated candidate
structures that chemistry-informed filters either propagated or
discarded. In each case, the model on its own would have produced
plausible-looking output that did not work. The system around the
model was what produced output that did.

A great deal of the book's content amounts to a careful unpacking
of this picture. What tools to give an agent, and how. How to
manage the context that the model sees on each turn. How to keep
the agent's behavior reliable as the surface of the task grows.
How to defend the system against inputs the model was never
trained to handle adversarially. None of this is glamorous;
all of it is decisive.

\section{The Approach}
\label{sec:intro-approach}

Three commitments shape the chapters that follow. None is
contentious in isolation; together they distinguish this book from
much of the surrounding literature, and the reader will benefit
from having met them before encountering the technical material.

\paragraph{Tools, not internals.}
Andrew Ng's framing of AI as ``the new electricity'' is by now a
well-worn analogy, and Karpathy has refined it specifically for
language models~\citep{karpathy2025software}: the model is a
utility, paid for by the token, served at scale by a small number
of providers, and consumed by everyone downstream. We adopt that
framing throughout. The model is treated as a service with a
calling convention, not as a research artifact whose internals we
are responsible for explaining. Where understanding \emph{how} a
transformer is trained would help an engineer make a better
decision, we mention the relevant fact and move on; where it
would not, we leave it out. The reader who wants the internals
has the original
literature~\citep{vaswani2017attention}, several
excellent textbooks, and a flourishing ecosystem of explainers;
we do not compete with them. The bargain is the same one
electrical engineers struck a hundred years ago: we do not need
to know how a generating station works in order to wire a
building, but we do need a precise mental model of voltage,
current, and the circuit at hand.

\paragraph{Reproducible claims.}
A great deal of writing about language models is, at the time of
writing, somewhere between vague advice and pseudoscience.
Recommendations of the form ``use this prompting trick to get
better results'' circulate freely without ever showing what
``better'' means or what the comparison is. We try not to do that
here. Every substantive claim about model behavior is paired
with a concrete, runnable example: a script in
\fname{manuscript/scripts/}, an output the reader can reproduce,
and where the result is sensitive to model choice, an explicit
note about what was used. The reader does not have to take our
word for any claim about prompting, tool calling, or evaluation;
the experiments are there to be run.

This commitment has an immediate consequence for the choice of
working model. Throughout the book, the default is Gemma~4 served
locally via Ollama. The reasons are practical, not ideological.
Closed-weight hosted models change underneath their identifiers
--- silently, often, and without notice. A reader who runs a
script today against \code{claude-3.5-sonnet} may, six months
from now, run the same script against a model that no longer
exists, or that returns different answers because the underlying
weights were updated, or because some internal knob was retuned.
Reproducibility against a moving target is, by definition, not
reproducibility. Open-weight models pinned to a specific release
do not have this problem: the artifact the reader downloads today
is the same artifact the book was written against, and the same
artifact a reader running the book in 2030 will be able to run.

\begin{updatebox}[title={Update --- the moving frontier}]
\small
Two ongoing developments illustrate why reproducibility against
closed-weight, hosted models is hard. In February~2026, telemetry
from thousands of Claude Code sessions showed a measured 67\%
drop in reasoning depth coinciding with Anthropic silently
lowering the default ``adaptive thinking'' level from high to
medium without any change to the model
identifier~\citep{adaptivethinking2026}. Around the same time,
Anthropic's Sonnet~3.5 and~3.7 --- workhorse models for many
production deployments only a year earlier --- entered late-stage
phase-out, with deprecation windows of weeks rather than years.
Neither development is malicious; both are inherent to the
closed-weight, hosted-service deployment model. Both are fatal to
reproducibility against a fixed identifier. A book whose
experiments are pinned to such a model would already, only a few
months after publication, no longer be runnable. A book whose
experiments are pinned to a downloadable release of an
open-weight model is.
\end{updatebox}

\paragraph{Learn by example, in pure Python.}
The third commitment is implementation discipline: the worked
code in this book is plain Python, written from first principles,
not built on agent frameworks. We use \code{requests} rather than
the OpenAI SDK; we reimplement a small wrapper around Ollama's
HTTP API rather than reach for \code{litellm}; we write our own
loop in Chapter~\ref{ch:toolcalls} rather than import a
\code{LangGraph} runtime. This is not a manifesto against
frameworks --- they are useful, especially for production systems
with operational concerns the book deliberately keeps out of
scope. It is a pedagogical choice. A reader who has built a
tool-calling loop with \code{requests.post} and a \code{for}~loop
has a working mental model that no amount of framework
documentation can substitute for. Once that mental model exists,
picking up a framework is straightforward; the reverse path is
much harder.

We do use general-purpose libraries where they earn their keep:
\code{pandas} for tabular work, \code{matplotlib} for plots,
\code{numpy} where vector operations matter. Nothing here is
reinvented for its own sake. The line we hold is at frameworks
that abstract away the agent loop itself, because that loop is
the subject matter.

Not every example is reproduced inline in the manuscript. Where
the example is short and central, it is in the text. Where it is
longer, the manuscript shows the salient pieces and points the
reader at the full source under \fname{manuscript/scripts/}.

\section{How to Read This Book}
\label{sec:intro-roadmap}

The book is divided into three parts. The first
(Chapters~\ref{ch:mental} and~\ref{ch:working})
develops a working understanding of what a language model is and
how to call one, at the level of detail a systems engineer needs.
The second (Chapters~\ref{ch:tools} through~\ref{ch:toolcalls},
plus the agents chapters that follow) covers the construction of
agents themselves: prompting craft, tool calling, control loops,
context engineering, multi-agent coordination. The third part
(Chapter~\ref{ch:reliable} and the closing chapters)
addresses the operational concerns of deploying and trusting
agents in production: security, evaluation, observability,
and the long-running questions of how to keep a system that
includes a language model from drifting out from under its
operators.

The chapters can be read in order, and that is the recommended
path on a first pass. A reader with prior LLM experience can skim
Chapters~\ref{ch:mental} and~\ref{ch:working} and start in earnest
at Chapter~\ref{ch:tools}. A reader whose interest is primarily in
production reliability will benefit from reading the early
chapters in any case --- the failure modes treated in
Chapter~\ref{ch:reliable} are easier to internalize after one has
seen the architectural patterns they apply to.

We will refer back to this introduction often. The mental picture
of \emph{model in a loop with tools} is the one the rest of the
book is built on.
